\section{Tables S1 and S2, with the observations identified}
\label{app:obs}

The two tables below are Tables~S1 and~S2 of the manuscript, with one column added at the
left. A label R\emph{x} marks a result this note discusses; a blank means the result is a
null that the manuscript never mentions, and there is nothing to say about it. Values are
the true ones: where they differ from those published, the difference is flagged and
explained in the note attached to the table.

\paragraph{On the window names.} The four windows are locked to the four trial events and
do not coincide with the intervals they are named after. The onset-locked window runs from
stimulus onset to two seconds, so after a $0.3$~s stimulus it contains $0.3$~s of stimulus
and $1.7$~s of delay. That is not a detail --- it is why the duration effects in that
window separate the two models --- and it is why this note calls the windows
\emph{onset-locked}, \emph{offset-locked}, \emph{signal-locked} and
\emph{response-locked} rather than the stimulus, delay, RS and response periods. We
suggest the manuscript do the same.

\subsection*{Table S1. Accuracy, response time, premotor and motor execution time}

\begin{center}\small
\begin{tabular}{llrrrrr}
\toprule
& effect tested & df1 & df2 & $F$ & $p$ & $\eta_p^2$\\
\midrule
\multicolumn{7}{l}{\emph{Accuracy}}\\
R1 & duration & 1 & 21 & 22.68 & $<$.001 & 0.519\\
R2 & coherence & 1.31 & 27.4 & 222.22 & $<$.001 & 0.914\\
R3 & duration $\times$ coherence & 2 & 42 & 12.17 & $<$.001 & 0.367\\
\multicolumn{7}{l}{\emph{Mean response time}}\\
R4 & duration & 1 & 21 & 65.79 & $<$.001 & 0.758\\
R5 & coherence & 1.50 & 31.39 & 14.19 & $<$.001 & 0.403\\
R6 & duration $\times$ coherence & 2 & 42 & 4.61 & .016 & 0.180\\
\multicolumn{7}{l}{\emph{Mean premotor execution time}}\\
R7 & duration & 1 & 21 & 69.75 & $<$.001 & 0.769\\
R8 & coherence & 1.59 & 33.35 & 14.11 & $<$.001 & 0.402\\
R9 & duration $\times$ coherence & 2 & 42 & 3.62 & .035 & 0.147\\
\multicolumn{7}{l}{\emph{Mean motor execution time}}\\
R10 & duration & 1 & 21 & 1.03 & .320 & ---\\
& coherence & 1.51 & 31.68 & 0.54 & .537 & ---\\
& duration $\times$ coherence & 2 & 42 & 2.39 & .104 & ---\\
\bottomrule
\end{tabular}
\end{center}

\noindent\emph{Note.} Values as published; this note did not recompute them, since its
own extraction differs from the manuscript's by about a hundred trials and substituting it
would confuse a difference of analysis with a difference of data. R10 labels the three
motor-execution-time rows together: the manuscript states them as one result, that no
factor affects motor execution time, and the note treats them as one.

\subsection*{Table S2. Rate of partial activations, responding channel}

\begin{center}\small
\begin{tabular}{llrrrrr}
\toprule
& effect tested & df1 & df2 & $F$ & $p$ & $\eta_p^2$\\
\midrule
\multicolumn{7}{l}{\emph{Onset-locked window, 0 to 2 s from stimulus onset}}\\
R11 & duration & 1 & 21 & 13.38 & .001 & 0.389\\
R12 & coherence & 1.33 & 28.01 & 5.25 & .021 & 0.200\\
R13 & time & 2.70 & 56.75 & 18.05 & $<$.001 & 0.462\\
& duration $\times$ coherence & 2 & 42 & 0.312 & .734 & ---\\
R14 & duration $\times$ time & 5.09 & 106.82 & 4.645 & $<$.001 & 0.181\\
R15 & coherence $\times$ time & 38 & 798 & 1.718 & .005 & 0.076\\
& duration $\times$ coherence $\times$ time & 38 & 798 & 1.058 & .377 & ---\\
\midrule
\multicolumn{7}{l}{\emph{Offset-locked window, 1 to 3 s from stimulus offset}}\\
R16 & duration & 1 & 21 & 33.19 & $<$.001 & 0.612\\
R17 & coherence & 2 & 42 & 3.36 & .044 & 0.138\\
R18 & time & 1.21 & 25.43 & 31.42 & $<$.001 & 0.599\\
R19 & duration $\times$ coherence & 2 & 42 & 3.474 & .040 & 0.142\\
R20 & duration $\times$ time & 5.38 & 112.99 & 4.520 & $<$.001 & 0.177\\
R21 & coherence $\times$ time & 38 & 798 & 0.737 & .879 & ---\\
& duration $\times$ coherence $\times$ time & 38 & 798 & 0.959 & .542 & ---\\
\midrule
\multicolumn{7}{l}{\emph{Signal-locked window, 0 to 0.2 s from the response signal}}\\
R22 & duration & 1 & 21 & 36.56 & $<$.001 & 0.635\\
R23 & coherence & 1.56 & 32.82 & 1.159 & .315 & ---\\
R24 & time & 1 & 21 & 21.96 & $<$.001 & 0.511\\
& duration $\times$ coherence & 2 & 42 & 0.041 & .959 & ---\\
R25 & duration $\times$ time & 1 & 21 & 7.046 & .015 & 0.251\\
& coherence $\times$ time & 2 & 42 & 0.407 & .668 & ---\\
& duration $\times$ coherence $\times$ time & 2 & 42 & 0.587 & .561 & ---\\
\midrule
\multicolumn{7}{l}{\emph{Response-locked window, $-0.2$ to 0 s before the press}}\\
R26 & duration & 1 & 21 & \textbf{7.613} & \textbf{.012} & \textbf{0.266}\\
& coherence & 1.90 & 39.98 & \textbf{1.845} & \textbf{.173} & ---\\
& duration $\times$ coherence & 1.86 & 39.10 & \textbf{0.342} & \textbf{.697} & ---\\
\bottomrule
\end{tabular}
\end{center}

\noindent\emph{Note.} Values reproduce those published exactly, to the last reported
digit, except in the response-locked window, where they are printed in bold because they
differ. Two corrected data files were supplied after the note was begun. The
signal-locked window did not reproduce with the earlier file and does with the corrected
one; nothing there changes. The response-locked window changes both in value and in
shape. In the corrected file the bin from $-0.1$ to $0$~s is exactly zero for every
participant and cell --- a partial activation precedes the response burst by definition,
and the burst begins about $160$~ms before the press --- so the window carries one
informative bin, the level and trend scores are the same number up to a constant, and the
four time-related rows of the published table are not separate tests. The three effects
above are what the window supports. The main effect of duration passes from $F=0.011$,
$p=.919$ as published to $F=7.613$, $p=.012$; Section~\ref{sec:corrected} sets out the
consequences.

\subsection*{One result outside the tables}

\noindent R27, the levelling off of the rate shortly before the response signal in the
two-second, fifty-per-cent condition. It is a reading of Figure~4A rather than a row of
any table, and the manuscript discusses it at length, proposing a strategic reduction of
temporal preparation to account for it. It is labelled here because
Section~\ref{sec:pred} derives it without that mechanism.
